\chapter{基于模型的数据分析}

\begin{quote}
仅仅因为我们手里有最好的锤子，并不意味着每个问题都是钉子。

--- 巴拉克·奥巴马
\end{quote}

\section{引言}

归根到底，本书描述的模型只有在能够回答科学问题时才真正有用。本章聚焦于如何用主动推断来理解经验数据。核心思想是：作为科学家，我们可以调用与前几章中假定大脑所使用的同一套数学工具。我们的总体目标，是恢复受试者大脑用来生成行为的生成模型参数，也就是其\textbf{主观模型}。为此，我们可以使用我们自己的生成模型，即“主观模型如何产生行为”的模型，也就是\textbf{客观模型}。我们可以根据观察到的行为，对客观模型进行反演，从而推断主观生成模型中的参数。这种元贝叶斯推断为我们提供了机会：一方面检验我们关于大脑所使用模型的假设，另一方面根据个体必须持有什么样的先验信念，才能使其行为成为贝叶斯最优，来对个体进行表型刻画。此类基于信念的计算表型分析，在计算精神病学、神经心理学和神经病学等新兴领域中具有重要前景。

\section{元贝叶斯方法}

本章讨论主动推断表述在行为实验数据分析中的用途。这已经超出了前几章中那些“原理验证”式的模拟，而是直接利用主动推断来回答科学问题。我们已经看到，在主动推断框架下，受试者的生成模型是决定行为的关键因素。这意味着，对于经验性行为测量之成因的假设，必须表述为：主体在选择这些动作时所使用的是哪一种备选生成模型。因此，我们面临的挑战，是通过调节生成模型的参数，也就是先验信念，将一个主动推断方案拟合到观察到的数据上。

广义上说，将计算模型拟合到观察到的行为上，有两个彼此相关的目的。第一个目的，是估计模型中最能解释某一特定受试者或某一组受试者行为的参数。这有助于用生成这些行为的计算过程来刻画主观行为，这一过程通常被称为\textbf{计算表型分析}。计算表型既可以与其他测量结合使用，例如建立神经影像结果与功能之间的联系，也可以单独使用，例如预测个体在其他情境下的行为表现，或预测治疗干预后的变化。

第二个目的，是比较不同的模型假设，也就是用不同模型来代表对某种行为现象的不同解释。这两项工作，即参数估计和模型比较，都对应于贝叶斯定理的同一侧。参数估计是在某个模型之下寻找参数设定的后验概率；模型比较则依赖于找到每个模型的边缘似然，也就是模型证据。回顾一下，贝叶斯定理为

\begin{equation}
P(u \mid \theta, m)\,P(\theta \mid m)=P(\theta \mid u,m)\,P(u \mid m).
\label{eq:chapter9-bayes}
\end{equation}

右侧对应的是：在给定模型 $m$ 和行为数据 $u$ 的条件下，参数 $\theta$ 的后验概率，以及模型证据；左侧告诉我们在建模时必须指定什么：我们需要对目标参数给出先验信念，并给出一个似然函数。

需要强调的是，虽然我们这里调用的仍是前几章中同样的贝叶斯推断方案，但其用途已经不同。原因在于，此处实际上同时存在两个推断过程。第一个过程是：生物体利用其关于感觉数据生成过程的模型，对世界状态以及该如何行动作出推断。这正是前几章的焦点。第二个过程是：作为客观科学家的我们，观察该生物体的行为，并试图通过反演我们自己的客观生成模型，来推断它正在使用的主观生成模型。换言之，我们此时是在对一个“推断过程本身”进行推断，这通常被称为\textbf{元贝叶斯}推断。

更形式化地说，这一路径将似然分布定义为一个主动推断问题的解。给定某一组参数，我们可以在主动推断之下模拟行为，并量化某一串动作会被采取的可能性。再结合我们对这些参数取值的先验信念，我们就得到了一个关于“生物体如何利用其生成模型来产生动作”的生成模型。虽然本章聚焦于主动推断，尤其是离散时间模型，但这里采用的一般方法也可以和任意似然函数配合使用。其他规范性行为模型，例如强化学习中使用的模型，也可以替代主动推断模型而被纳入这一框架。

接下来的部分将展开一个可用于元贝叶斯推断的通用推断方案示例，即变分拉普拉斯，以及用于模型比较的层级模型。之后，我们会给出一套基于模型的数据分析简明操作流程，并回顾一个关键应用案例。需要强调的是，要有效使用这些方法，并不要求读者完全掌握其中的技术细节；因此，如果读者对这些细节兴趣不大，可以跳过第 9.3 节和第 9.4 节。

简言之，基本思想是：评估在给定未知目标参数，也就是受试者先验信念参数时，任何一组被观察到的选择序列出现的似然。然后，把这一似然与我们对这些参数的客观先验结合起来，按通常方式求得受试者先验的后验分布。如果我们拥有多个受试者，就可以使用参数经验贝叶斯（PEB）把这些后验组合起来，从而对群体效应或被试间差异作出推断。这里所需的似然，本质上就是：在受试者对动作的后验信念之下，采样到该观察到的选择序列的概率。这些后验信念取决于受试者看到了什么，也就是线索或刺激，以及她原本持有的先验信念；而它们可以通过求解相应的主动推断方案来直接得到。要注意，我们在这里实际上使用了两次贝叶斯程序：第一次是为了求出受试者关于动作的后验信念；第二次是为了求出我们关于那些刻画受试者的未知先验参数的后验信念。下面我们依次梳理这一元贝叶斯程序的各个组成部分。

\section{变分拉普拉斯}

变分拉普拉斯是一种与预测编码建立在同样原理之上的推断方案。不过，它可以用于比前面所遇到的高斯型似然更一般的似然函数。本节先给出这里关心的似然函数 $L(\theta)$ 的概览。该似然函数应当给出：当生成模型的参数取值为 $\theta$ 时，一个主动推断方案产生相应动作的概率。被选中的动作依赖于所获得的观测：

\begin{equation}
\begin{aligned}
L(\theta) &= \ln P(\tilde{u}\mid \theta, m,\tilde{o}), \\
P(\tilde{u}\mid \theta,m,\tilde{o}) &= \tilde{u}\cdot \sigma\!\bigl(\theta_{\alpha}\ln \mathbf{u}\bigr), \\
\mathbf{u} &= \pi \cdot U, \\
\pi &= \arg\min_{\pi} F.
\end{aligned}
\label{eq:chapter9-likelihood}
\end{equation}

把它拆开来看，第一行给出的是：观察到的动作序列 $\tilde{u}$ 的对数似然，它是参数 $\theta$、模型 $m$ 以及真实实验中呈现的刺激序列 $\tilde{o}$ 的函数。动作序列出现的概率，是通过把这些参数写入某个第 7 章那样的 POMDP 模型中的先验信念来计算的。然后，我们像第 4 章和第 7 章那样去求解该 POMDP，不过在模拟中强制它采取已经观察到的动作序列，并向它呈现与真实实验相同的刺激。正如前几章所述，这一过程会得到一个合成受试者关于其所选择策略或行动路径的信念 $\pi$。这些信念使其世界生成模型的自由能 $F$ 达到最小。接着，我们就可以利用这些策略信念，计算执行某一动作序列的平均概率。为此，需要把每个策略的概率分配到该策略所隐含的动作上，动作由数组 $U$ 索引。最后，引入一个 softmax 温度参数 $\theta_{\alpha}$，以刻画模型未能解释的行为随机性，也就是“手抖性”。如果这个 softmax 参数等于 $1$，我们实际上是在假设受试者按照其关于动作的后验信念进行采样，这有时也称为\textbf{匹配行为}。相反，如果这个 softmax 参数非常大，那么实际发出的动作就是主观后验概率最大的那个动作，也就是说，受试者总是选择最可能的选项。这个 softmax 参数本身也可以被估计。

由此，我们得到了在给定刺激序列与参数的条件下，模型产生这些动作的概率，也就是“在某个模型下，行为数据的似然”。如果进一步为客观层参数赋予高斯先验
\[
\theta \sim \mathcal{N}\!\bigl(\eta,\Pi^{(1)}\bigr),
\]
那么我们就可以借助拉普拉斯假设，写出对模型证据的一个自由能近似：

\begin{equation}
\begin{aligned}
\ln P(\tilde{u}\mid m,\tilde{o})
&\approx L(\mu)+\frac{1}{2}\,\varepsilon^{\mathsf T}\Pi^{(1)}\varepsilon
+\ln \left|\nabla_{\mu\mu}L(\mu)-\Pi^{(1)}\right|, \\
\varepsilon &= \eta-\mu, \\
\mu &= \arg\max_{\mu}\left\{
L(\mu)-\frac{1}{2}\left(
\varepsilon^{\mathsf T}\Pi^{(1)}\varepsilon
+\ln \left|\nabla_{\mu\mu}L(\mu)-\Pi^{(1)}\right|
\right)\right\}.
\end{aligned}
\label{eq:chapter9-evidence}
\end{equation}

式 \eqref{eq:chapter9-evidence} 与框 4.3 中展开的结果本质相同，只是这里推广到了多维参数空间，并显式代入了后验协方差的形式，同时假定先验服从正态分布。在第 4 章以及第 8 章中的应用里，我们忽略了那些不依赖于众数的项；但当我们讨论模型比较问题时，这些项就必须保留。

为了找到使式 \eqref{eq:chapter9-evidence} 最后一行达到最大值的 $\mu$，我们可以执行梯度上升。在二次近似下，这可化为

\begin{equation}
\dot{\mu}=\nabla_{\mu}L(\mu)+\Pi^{(1)}\varepsilon.
\label{eq:chapter9-gradient}
\end{equation}

虽然这里使用的对数似然梯度未必总能写出显式形式，但我们可以利用有限差分方法对其进行合理的数值近似。同样的方法也可用于求取后验精度，也就是负对数似然的二阶导数或 Hessian，再加上先验精度。式 \eqref{eq:chapter9-gradient} 是最简单的更新形式，但在实践中往往会使用基于局部曲率的更复杂方法。

\section{参数经验贝叶斯（PEB）}

上一节中的变分拉普拉斯程序，使我们能够对选择行为模型进行推断，并量化该模型的证据。这让我们可以对个体进行计算表型分析，也可以比较关于该个体的不同解释性假设。然而，更有趣的问题往往发生在群体层面。举例来说，我们可能关心某个参数，例如先验偏好的精度，如何随着年龄变化。为了回答这个问题，我们可以先使用第 9.3 节的方法，对不同年龄个体的行为分别拟合模型；然后再建立一个生成该参数的广义线性模型，并把年龄纳入其中：

\begin{equation}
P(\theta \mid \beta,X)=\mathcal{N}\!\bigl(X\beta,\Pi^{(2)}\bigr).
\label{eq:chapter9-peb}
\end{equation}

这里，$X$ 是一个矩阵，它的列表示不同的解释变量，它的行对应不同参与者。通常，$X$ 的第一列是一列全 $1$，用来表示受试者总体平均参数的效应；在本例中，第二列则可以是每位参与者的年龄。向量 $\beta$ 表示 $X$ 中每个解释变量的效应大小。于是，$\beta$ 的第一个元素就是该精度参数的平均值，第二个元素则表示年龄对精度的影响。这个值也就是在“年龄为横轴、预测精度为纵轴”的图中那条直线的斜率。$X$ 可以有任意多列，$\beta$ 也可以有任意多个元素。

一旦拟合了式 \eqref{eq:chapter9-peb} 所表达的模型，并为 $\beta$ 设定先验，我们就可以询问各个解释变量的作用。例如，我们可以通过比较两个模型的证据，来判断年龄是否影响先验偏好的精度：一个模型允许 $\beta$ 的第二个元素偏离零，另一个模型则对它是零持有精确先验。实践中，这通常可以借助贝叶斯模型约简来完成，而不必重新反演多个模型。

\section{基于模型分析的操作说明}

在实践中，我们通常按照下面的步骤，使用主动推断来分析经验性选择行为。这些步骤对应于 SPM12 Matlab 包中的相关程序。

\begin{enumerate}
\item 收集行为数据，包括个体做出的选择，以及该个体在作出该选择时可获得的感觉输入。此外，还要收集用于第二层被试间分析的数据，例如受试者是否属于患者组或对照组、相关人口统计信息等。
\item 像第 7 章那样构造一个 POMDP 模型。它应当是一个以参数为输入、并输出一个已被完整指定但尚未求解的 POMDP 的函数。
\item 指定一个似然函数，也就是式 \eqref{eq:chapter9-likelihood}。它规定模型应如何被用来计算似然。通常，这一步会调用一个 POMDP 求解器，例如 \texttt{spm\_MDP\_VB\_X.m}，来模拟行为，并量化观察到的动作之似然。
\item 用期望与精度来指定参数的先验信念。通常这些先验会以 $0$ 为中心，而精度则反映参数可能取值的合理范围。
\item 求解后验概率和模型证据。这里通常使用一种标准推断方案，例如上面介绍的变分拉普拉斯程序，也就是式 \eqref{eq:chapter9-gradient} 所对应的方法。例行程序 \texttt{spm\_nlsi\_Newton.m} 可以自动完成这一过程。
\item 进行群体层分析。通常使用 PEB，把每个个体的已估计参数视为由第二层模型生成。这样就可以检验关于这些参数成因的假设。实践中，这一步可使用 \texttt{spm\_dcm\_peb.m}。其他可选分析还包括：把每位受试者的推断参数与其他个体特异性测量之间做标准统计关联分析。例如，可以使用典型变量分析来评估问卷分数与推断参数之间的关系。
\end{enumerate}

图 9.2 对上述步骤的概括，是以第 7 章所描述的 T 迷宫任务中的老鼠行为为例的。首先，我们把一只老鼠放入 T 迷宫，在左臂或右臂之一放置奖励性刺激，并在中央臂放置信息性线索。然后记录老鼠采取的动作序列。这一程序可以在多个试次中重复，以记录学习后的行为；也可以在多只不同老鼠、不同干预条件下重复，例如药理学干预或光遗传学干预。

一旦获得这些行为数据，我们就需要一个似然函数，用来量化：在特定参数设定之下，某一只老鼠在某一条件下产生这些行为的概率。为此，我们可以构造第 7 章讨论过的 POMDP 模型。这个模型必须按我们希望求似然的参数来参数化。例如，如果我们想评估偏好相关的精度，就可以引入一个对数尺度参数，使偏好分布变得更尖锐或更平坦。

在搭建好从老鼠视角出发的生成模型后，我们就能利用第 4 章中的信念更新方程，自动求解该 POMDP。这使得我们能够计算：在某一特定偏好尺度参数取值下，给定模型产生这些数据，也就是访问迷宫不同臂的序列的概率。再把这一似然与我们的先验结合起来，就完成了从科学家视角出发的行为生成模型的指定。然后，就可以用变分拉普拉斯，为每只老鼠求出该尺度参数的后验分布。

在真正分析实际数据之前，一个合理步骤，是先用该 POMDP 模型生成虚构数据，检查其表面有效性，看看这些数据在定性上是否与手头问题相符。另一个明智测试是\textbf{参数恢复}：在某些已知参数设定之下生成虚构数据，再反演模型，看看这些参数能否被恢复出来。这有助于验证某些参数，或者它们的组合，是否是可恢复的，也就是是否可识别。

最后，我们可以为一个线性模型构造设计矩阵，其中每一行代表一种计算表型，例如每位受试者偏好参数的后验密度，而各列则代表这些受试者的不同属性。这些属性可能解释不同老鼠偏好的差异。除了一列全 $1$，用于表示所有老鼠的平均偏好，还可以加入年龄、是否接受药物、等等变量。借助这个被试间效应模型，我们就可以进一步进行 PEB 分析，以评估这些解释变量对先验偏好的贡献。

\section{生成模型示例}

本节借助文献中的两个例子，说明如何使用连续型和离散型生成模型。首先，我们简要回顾 Adams、Aponte 等人（2015）以及 Adams、Bauer 等人（2016）的做法。下文统称其为 Adams 等人。他们用平滑追踪眼动任务来量化每位受试者生成模型中的精度参数。该研究设计的重要之处在于：它同步采集了电生理数据，也就是脑磁图数据，使作者能够进一步询问精度或置信度编码的神经生物学基础。接着，我们再看 Mirza 等人（2018）的研究，他们把扫视眼动建模为一个 POMDP，用以分析主动信息搜寻行为。图 9.3 以示意方式给出了这两个实验。希望这两个例子能帮助读者理解：上文概括的一般方法，是怎样被真正用于回答科学问题的。

如果按照图 9.2 中的步骤来看，Adams 等人首先进行了数据收集，也就是步骤 1。他们让受试者盯住一个移动的视觉目标，并维持注视。具体任务细节并不重要，但该任务有两个条件：在第一种条件下，目标沿可预测的正弦轨迹移动；在第二种条件下，目标沿相同轨迹移动，但叠加了高斯噪声。收集到的数据包括眼动轨迹。接着，作者建立了一个主观模型，也就是步骤 2。与图 9.2 中的 POMDP 不同，他们选择的是第 8 章那类连续时间模型。简言之，该模型预测来自眼睛的本体感觉输入和视觉输入，并假设注视点会被目标位置所吸引。步骤 3 中的似然，是借助第 4 章概述的主动预测编码方案构造出来的。这个似然量化的是：在给定一组对数精度参数以及步骤 2 中所设模型时，观测到这些眼球运动动作的概率。再往下的步骤 4，作者为这些对数精度指定了正态分布先验。随后，他们进行模型反演，也就是步骤 5，从而得到这些精度参数的后验分布。至于步骤 6，他们没有使用 PEB，而是把与行为任务同步采集的神经影像数据纳入分析。他们利用动态因果建模估计了初级视觉皮层浅层锥体细胞的增益。这样一来，他们就为每个受试者同时得到了精度参数估计和突触增益估计，从而可以在群体层面检验主观模型参数及其生物学基底之间的相关关系。这个相关性的展示，提供了一个重要例子：主动推断形式的行为模型如何使我们能够提出并回答“信念更新与神经生物学之间关系”的问题。

第二个例子中，Mirza 等人使用 POMDP 形式的主动推断，研究信息增益在驱动人类行为中的作用。仍按图 9.2 的步骤来看，他们首先收集行为数据，也就是步骤 1。当时，受试者执行的是一个视觉觅食任务：目标是把一个视觉场景归入若干类别之一。场景的每个部分，只有当受试者把目光落在相应位置时才会显露出来，因此受试者必须进行多次注视，才能获得足够证据来判断场景类别。作者收集的数据包括扫视，也就是快速眼动的序列。步骤 2 中使用的模型，是 Mirza 等人（2016）描述过的一个 POMDP 模型。它根据当前注视位置和场景类别，来预测离散化的本体感觉、视觉和反馈结果。偏好，也就是第 7 章中的 $C$，被施加在反馈结果上，使模型预期并因此偏好自己分类正确。步骤 3 中的似然函数，是在不同参数设定下用第 4 章中的方案求解该模型而得到的。这里涉及的参数包括一个偏好分布精度的对数尺度参数，以及其他一些参数。步骤 4 中，作者为该对数尺度和其他参数指定先验分布，并在步骤 5 中为每一位受试者分别反演模型。他们利用每位受试者的对数证据，比较了两类模型：一类把期望自由能中的认知成分纳入行为动机，一类则没有。结果发现，对所有受试者而言，包含认知可供性的模型都具有更高证据。随后，他们在步骤 6 中使用 PEB，评估受试者在多次试验过程中先验信念如何变化，结果发现证据支持信念参数发生了变化，也就是出现了主动学习。最后，他们又使用典型协变量分析，来评估两类量之间的关系：一类是每位受试者的表型变量线性组合，例如偏好精度；另一类是其行为表现指标的线性组合，例如正确率和反应时。

\section{错误推断模型}

\begin{quote}
人们通常都会很快相信那些他们希望为真的东西。

--- 尤利乌斯·凯撒
\end{quote}

鉴于这些方法对计算精神病学等领域的重要性，本章最后概述一下\textbf{错误推断}。这一主题之所以关键，是因为精神病理学常常可以被理解为信念更新失败的结果。像主动推断这样的推断框架有一个重要优势：它能够同时处理精神障碍的多个层面，把适应不良行为，例如强迫、成瘾，与心理层面现象，例如错误信念，以及生物层面现象，例如神经调质异常，连接在同一个解释框架中。

由于这里不可能公正覆盖利用主动推断建模疾病过程的庞大文献，本节只提供一个高度压缩的概览，目的是给出一个思考计算病理学的框架。表 9.1 给出了一组并不穷尽的示例，它们既包括离散时间模型，也包括连续时间模型，分别可追溯到第 7 章和第 8 章中的方案。下面的讨论主要借助 POMDP 一类模型的结构来展开；而支撑这些情境中错误推断的原则，其实大体相同。

这类推断路径，也就是主动推断等方法，背后的基本假设是：精神病理状态可以被概念化为\textbf{推断障碍}。这里所谓“障碍”并不一定意味着推断机制本身是坏的，例如后验概率计算错误。事实上，在表 9.1 所回顾的大多数研究中，推断机制本身运行正常，只是它所依据的生成模型存在缺陷，也就是生成模型中被赋予了异常的先验信念。因此，从根本上说，病理性表现是异常先验信念的后果；而这些先验是可以通过本章介绍的基于模型的数据分析方法恢复出来的。

下面按照若干病理类型做简要概括。

\subsection*{成瘾、冲动性与强迫性}

成瘾是一个重要例子：表面上看，它是一类异常行为，但在合适的生成模型之下，它同样可以被表述为最优推断。相关研究利用“有限报价任务”等范式说明了这一点：在任务中，参与者对“等待后是否会获得奖励”的把握程度不同；较低的置信度会导向与成瘾相关的强迫性行为。后续工作还进一步分析了更冲动或更不冲动的行为背后对应的先验信念，并在“离开资源斑块”范式中考察了强迫症里先验精度减弱所扮演的角色。

\subsection*{妄想}

妄想以固定而错误的信念为特征。在主动推断中，它可以被直接刻画为：在缺乏支持性证据时，后验概率分布却异常精确。如果这种后验足够精确，那么即使后来出现相互矛盾的证据，它也依旧会保持不变。不同妄想背后的机制可能并不相同。例如，感觉衰减失败可能是“行动归属妄想”的关键原因。近期研究还提供了共享性妄想的例子，即两个没有信息的主体如何就世界状态达成高度自信的一致判断。

\subsection*{幻觉}

这类模拟通常建立在先验精度与似然精度失衡之上。若似然精度的衰减不足，就会过度解读偶然的感觉数据；若先验精度的衰减过强，又会使先验信念无法被经验修正。这两条路径都可能构成错误知觉推断的机制，也就是幻觉的机制。

\subsection*{人际与人格障碍}

人际推断依赖于我们对他人以及他们会如何回应自己决策的模型。这推动了信任博弈模型以及慈善博弈模型的发展。后者被用于复现心理病态中常见的自我膨胀与缺乏悔意等特征。相关特质可以通过调节如下因素来模拟：关于自我价值的信念在多大程度上依赖于“慈善”还是“自利”的选择，以及主体对他人认可的敏感程度。

\subsection*{眼动综合征}

在一些论文中，连续时间生成模型被用来预测牛顿系统的动态演化。通过让生成模型的某些部分在条件上与其他部分独立，就可以诱发某些眼动综合征，例如核间性眼肌麻痹。

\subsection*{药物治疗}

结合第 5 章中关于精度参数与神经化学物质之间关系的讨论，可以设想通过模拟药理操控来研究这些系统。相关工作展示了若干合成药理干预对一个眼动延迟任务表现的影响，从而证明：这些方法不仅可以模拟病理状态，也可以模拟治疗措施的作用。

\subsection*{前额叶综合征}

一些模拟通过削弱转移精度，区分了内侧前额叶与外侧前额叶综合征的差异：前者更影响任务动机，后者更影响记忆引导任务的表现。

\subsection*{视觉忽略}

对空间左侧的忽视，可以通过多种不同的“病损先验”诱发。例如，若左侧空间的狄利克雷参数升高，则朝左扫视所带来的新奇性就会降低，导致个体更多去采样右侧。又或者，如果把偏好设置为与右侧本体感觉或视觉结果一致，或者增强向右扫视的习惯性倾向，也会重现定性上类似的行为。

\subsection*{内感受推断障碍}

关于内感受推断，也就是主动推断在内感受领域中的应用，相关模拟表明：例如关于心脏或胃部信号的先验精度与似然精度一旦失衡，就会导致关于身体内部状态的错误信念、对身体症状的错误知觉，甚至心身性幻觉。此外，这种失衡还会级联性地影响自主调节与行动选择，进而引发多种适应不良行为，例如过度警觉、过度用药以及过度节食。

异常先验既可以是关于状态或精度的先验，也可以是关于生成模型形式本身的结构先验。要理解病理行为的成因，一个有用的思路，是去分析用于策略的先验信念，并考察其中每一部分的破坏会如何导致异常的策略选择。策略先验依赖于期望自由能，而期望自由能又依赖于后验信念、信息增益的潜力以及先验偏好 $C$。此外，策略先验中还可以加入一个固定形式项 $E$，用来表示习惯性偏置。

依次来看。首先，后验信念依赖于先验和似然。若要形成异常的后验信念，则其中一者或二者必定出了问题。最常见的情况，是精度平衡被高估或低估。当似然精度相对于先验精度过高时，主体就会过度解读感觉输入，即使这些感觉输入本身带有噪声；这会导致一种“过拟合”式的状态，也就是从偶然数据中得出不恰当结论。若精度平衡朝相反方向偏移，即过度偏向先验置信度，那么由内部生成的知觉就会对相矛盾的感觉输入表现出顽固抵抗。这两种机制都与幻觉的形成有关；在层级模型中，它们甚至可以并存。考虑到多种精度与神经调质系统之间的联系，这也使得一些伴有神经化学异常的疾病，例如胆碱能系统受损的路易体痴呆，以及多巴胺系统异常的精神分裂症，出现幻觉现象就显得合理了。

其次，我们再看信息增益的作用。在这里，似然精度与先验信念精度共同决定了：不确定性在多大程度上是可以被消解的，以及当前究竟有多少不确定性需要被消解。先验精度既可以作用于生成模型参数，也就是影响新奇性，也可以作用于状态，也就是影响显著性。若把条件概率参数理解为突触效能，或者把精度理解为突触增益，那么突触断连综合征就可以被看作这二者之一，或二者同时，发生了破坏。失去的突触无法被调制，因此它非常类似于对某个条件概率持有极其自信的先验，因为新数据已无法再更新对应的效能。这对不同策略所能带来的信息增益有重要影响，而这正是感官忽略综合征建模中被利用的机制。

最后，偏好与策略先验对行为具有清晰而直接的影响。它们可能构成成瘾性习惯的基础，也可能构成许多精神病学和神经病学综合征中冷漠状态的基础。总之，在上述生成模型的不同位置上出现的有缺陷先验信念，可以为病理行为提供一种功能性或目的论层面的解释。

\section{总结}

本章概述了一种把前几章中理论模型用于经验数据发问的方法。借助这种方法，主动推断可以作为一种非侵入式工具，用来探查个体在决策时所依赖的计算过程。我们聚焦于若干相对简单的例子，但实际上，基于主动推断的模型已经被扩展到更真实、更复杂的任务上，以便诱发更丰富的行为，从而支持计算表型分析。

除了给出一个六步式的基于模型分析流程外，我们还强调了两个具体应用案例。这两个案例突出了这一流程在实践中的若干关键变化维度，包括所测量行为的类型，例如连续轨迹还是离散选择，所使用模型的类型，例如连续模型还是离散模型，以及最关键的，研究者真正想回答的科学问题。正是最后这一点决定了前面各项建模选择。我们看到，计算表型分析既可以和神经影像结合起来，用以探究突触增益与精度之间的关系；也可以通过模型反演来评估不同行为驱动力及其绩效预测因子的贡献。归根到底，图 9.2 中的六个步骤提供了一种通用方法，用来设计实验，以非侵入方式审问人类或其他动物用来驱动行为的隐含生成模型。这为我们回答“神经系统在健康与疾病中的功能如何运作”提供了现实路径。
